% ================================================================
% ==== FILE: content/k_levels/k5.tex
% ================================================================

% ==============================
%  Ontology of Continua — Core
%  K-level Module: K5
%  File: content/k_levels/k5.tex
%  Status: FULLY DEFINED
% ==============================



\paragraph{Canonical tuple projection note.}
Local K-level displays in this module are projections of the canonical public tuple \(K=(\Omega,\partial\Omega,A,\Theta,P,J,C,k,M)\). When a display omits \(M\), reorders components, or uses level-indexed shorthand, the omitted embedding context is inherited from the surrounding level discussion rather than introduced as a competing definition.

\subsubsection{\texorpdfstring{$K_5$}{K_5} Übersicht}
\label{sec:k5-overview}

$K_5$ ist das Niveau, bei dem \emph{elektrische Erregbarkeit} als a auftritt
neue ontologische Dimension. Die Protozellengrenze von $K_4$ wird zu einer
\emph{elektrisch aktive Membran} durch die Bildung von Ionenkanälen
mit kontrollierter Leitfähigkeit. Dadurch wird eine neue Achse generiert:
\[
A_{\mathrm{exc}}: \Delta V \;\mapsto\; \text{elektrische Zustandsdifferenz},
\]
wobei $\Delta V$ das Transmembranpotential ist.

Diese Achse ist in $K_4$ nicht definierbar. Es stellt eine neue Klasse von dar
unvereinbare Unterschiede und stellt daher eine echte Steigerung dar
Dimension.

Organisatorisch zeichnet sich $K_5$ aus durch:
\begin{itemize}
    \item stabile Ionenkanäle (offen/geschlossen/undicht/blockiert),
    \item kontrollierte Ionenflüsse,
    \item erregbare Zyklen (Anregung–Erholung),
    \item Feuerfestdynamik,
    \item Lärmschutz und Schwellenreaktionen,
    \item autarke elektrische Organisation, eng gekoppelt an die
          chemischer Innenraum.
\end{itemize}

$K_5$ ist das minimale \emph{neuronenartige} Kontinuum, auch wenn es kein reales ist
Neuronen gibt es noch.

% ================================================================
\subsubsection{Zustandsraum \texorpdfstring{$\Omega(K_5)$}{\Omega(K_5)}}

Ein Zustand von $K_5$ umfasst:
\[
\Omega(K_5) = \{ \Delta V,\, c_i,\, g_{\mathrm{ion}},\,
                P_{\mathrm{open}},\,
                x(t),\,
                G=(V,E),\,
                I(t) \}.
\]

Komponenten:
\begin{itemize}
    \item $\Delta V$ – Transmembranpotential,
    \item $c_i$ — Ionenkonzentrationen innen und außen,
    \item $g_{\mathrm{ion}}$ — Leitfähigkeiten von Ionenkanälen,
    \item $P_{\mathrm{open}}$ – offene Wahrscheinlichkeiten,
    \item $x(t)$ – interne Aktivierungsvariablen,
    \item $G=(V,E)$ – protoneuronales Konnektivitätsdiagramm für mehrere Einheiten
          $K_5$ Systeme,
    \item $I(t)$ – externe oder interne Ströme.
\end{itemize}

Voraussetzung für die Zulassung sind:
\[
|\Delta V| < \Theta_{\mathrm{volt-max}}, \quad
c_i \in \text{Lebensfähigkeitsbereich}, \quad
0 \le P_{\mathrm{open}} \le 1.
\]

% ================================================================
\subsubsection{Grenze \texorpdfstring{$\partial\Omega(K_5)$}{\partial\Omega(K_5)}}

Die Membran erbt die Grenze von $K_4$ mit einer neuen Struktur:
\[
\partial\Omega(K_5)
    = \partial\Omega(K_4)
      \cup \{ \text{Ionenkanäle mit Gating-Kinetik} \}.
\]

Für jeden Kanal $k$:
\begin{itemize}
    \item Zustandsraum: offen/geschlossen/blockiert/undicht,
    \item-Parameter: $g_k$, Umkehrpotential $E_k$, Gating-Variablen,
    \item Lokalisierung auf Patch-Ebene,
    \item Wechselwirkungen mit lokaler Krümmung und Spannung.
\end{itemize}

Der Zusammenbruch der Kanalordnung oder der Gating-Kinetik zerstört $\Omega(K_5)$.

% ================================================================
\subsubsection{Achsen \texorpdfstring{$A(K_5)$}{A(K_5)}}

Die minimalen Achsen sind:
\[
A(K_5) = A(K_4)\; \Tasse\;
\{ A_{\mathrm{exc}},\; A_{\mathrm{ion}},\; A_{\mathrm{gate}} \}.
\]

Wo:
\begin{itemize}
    \item $A_{\mathrm{exc}}$: elektrische Erregerachse,
    \item $A_{\mathrm{ion}}$: Ionenspezies und Flussunterschiede,
    \item $A_{\mathrm{gate}}$: Gating-Zustandsunterschiede für Ionenkanäle.
\end{itemize}

Diese Achsen sind mit rein chemischen Differenzen von $K_4$ und nicht kompatibel
implizieren ein höherdimensionales Kontinuum.

% ================================================================
\subsubsection{Potentiale \texorpdfstring{$P(K_5)$}{P(K_5)}}

Zu den Potenzialen gehören:
\[
P(K_5)=\{
E_{\mathrm{ion}},\;
P_{\mathrm{gate}},\;
P_{\mathrm{exc}},\;
P_{\mathrm{Rauschen}},\;
P_{\mathrm{refrac}},\;
P_{\mathrm{Kopplung}}
\}.
\]

Interpretation:
\begin{itemize}
    \item $E_{\mathrm{ion}}$ — Nernstpotentiale für jede Ionenart,
    \item $P_{\mathrm{gate}}$ – Gating-Aktivierungspotentiale,
    \item $P_{\mathrm{exc}}$ — Anregungsschwellen für die Spike-Initiierung,
    \item $P_{\mathrm{noise}}$ — Lärmschutzpotentiale,
    \item $P_{\mathrm{refrac}}$ — Refraktär-Rückstellpotentiale,
    \item $P_{\mathrm{Kopplung}}$ — Potentiale für die elektrische Kopplung
          mehrere Einheiten über $G$.
\end{itemize}

% ================================================================
\subsubsection{Schwellenwerte \texorpdfstring{$\Theta(K_5)$}{\Theta(K_5)}}

Schwellenstruktur:
\[
\Theta(K_5)=\{
\Theta_{\mathrm{exc}},\;
\Theta_{\mathrm{Rauschen}},\;
\Theta_{\mathrm{volt-max}},\;
\Theta_{\mathrm{channel-open}},\;
\Theta_{\mathrm{channel-close}},\;
\Theta_{\mathrm{Kanalrauschen}},\;
\Theta_{\mathrm{front}},\;
\Theta_{\mathrm{refrac}}
\}.
\]

Bedeutung:
\begin{itemize}
    \item $\Theta_{\mathrm{exc}}$: Anregungsschwelle für Spike
          Generation,
    \item $\Theta_{\mathrm{noise}}$: Rauschgrenze, jenseits derer Anregung
          ist verloren,
    \item $\Theta_{\mathrm{volt-max}}$: maximal sichere Transmembran
          Spannung,
    \item $\Theta_{\mathrm{channel-open}}$, $\Theta_{\mathrm{channel-close}}$:
          Gating-Lebensfähigkeitsschwellen,
    \item $\Theta_{\mathrm{channel-noise}}$: Rauschtoleranz beim Gating,
    \item $\Theta_{\mathrm{front}}$: minimale Bedingungen für die Ausbreitung
          Fronten in Mehrfamilienhäusern,
    \item $\Theta_{\mathrm{refrac}}$: minimaler Refraktärerholungsgrad.
\end{itemize}

Überschreitung von $\Theta_{\mathrm{volt-max}}$ oder $\Theta_{\mathrm{channel-noise}}$
vernichtet das Kontinuum.

% ================================================================
\subsubsection{Flüsse \texorpdfstring{$J(K_5)$}{J(K_5)}}

Flüsse:
\begin{itemize}
    \item Ionenströme:
    \[
    J_{\mathrm{ion},k} = g_k(\Updelta V - E_k),
    \]
    \item Leckströme und Rauschströme,
    \item Gating-Variablenflüsse:
          $\dot{x} = f(x,\Delta V)$,
    \item Spike-Propagation fließt über Graphkanten $E$,
    \item Rückgewinnungsströme der Feuerfestigkeit,
    \item-Kopplungsflüsse zwischen Einheiten.
\end{itemize}

Stabilität erfordert:
\[
\sigma(J) < \sigma_{\max}(K_5).
\]

% ================================================================
\subsubsection{Zyklen \texorpdfstring{$C(K_5)$}{C(K_5)}}

Charakteristische Zyklen:
\[
C(K_5)
= \{ C_{\mathrm{exc}},\;
     C_{\mathrm{refrac}},\;
     C_{\mathrm{Lärmschutz}},\;
     C_{\mathrm{Kanaldynamik}},\;
     C_{\mathrm{Kopplung}} \}.
\]

Beschreibungen:
\begin{itemize}
    \item $C_{\mathrm{exc}}$: Anregungszyklus (Ruhe → Anregung →
          Peak → Repolarisation),
    \item $C_{\mathrm{refrac}}$: Kanal zur Wiederherstellung des Refraktärzyklus
          Staaten,
    \item $C_{\mathrm{noise-control}}$: Unterdrückung der Stochastik
          Schwankungen,
    \item $C_{\mathrm{channel-dynamics}}$: Gating-Zyklen, die aufrechterhalten
          Erregbarkeit,
    \item $C_{\mathrm{Kopplung}}$: Synchronisationszyklen mit mehreren Einheiten.
\end{itemize}

Jeder Zyklus hat eine charakteristische Periode $\tau_{C_j}$.

% ================================================================
\subsubsection{Zeit \texorpdfstring{$\tau(K_5)$}{\tau(K_5)}}

Die Zeit ergibt sich aus der Existenz stabiler Anregungszyklen:
\[
\tau(K_5)
    = \min_j \tau_{C_j},
\qquad
\Pi(C_j) > \Theta_{\mathrm{Zeit}}.
\]

Aufgrund der Schärfe ist die zeitliche Ordnung bei $K_5$ starrer als bei $K_4$
Erregungsschwellen und Refraktärintervalle.

% ================================================================
\subsubsection{Kontinuität \texorpdfstring{$k(K_5)$}{k(K_5)}}

Unter Verwendung der allgemeinen Definition:
\[
k_5
= H(\Omega(K_5))
  \cdot\frac{|V_{\mathrm{Zyklus}}|}{|V|}
  \cdot\frac{\sum_j C_j^{\mathrm{eff}}}{\sum_j C_j^{\max}}
  \cdot\left(1-\frac{\sigma(J)}{\sigma_{\max}}\right)
  \cdot\left(1-\frac{T(K_5)}{\Theta_{\mathrm{stab}}}\right)_+.
\]

Interpretation:
\begin{itemize}
    \item Zyklusfraktion steuert Lebensfähigkeit,
    \item Übermäßiger Lärm reduziert $k_5$,
    \item zu groß $\Delta V$ kollabiert $k_5$ auf Null,
    \item koordinierte Anregung erhöht $k_5$.
\end{itemize}

% ================================================================
\subsubsection{Strukturelle Spannung \texorpdfstring{$T(K_5)$}{T(K_5)}}

Spannungsquellen:
\begin{itemize}
    \item inkompatible Ionengradienten,
    \item Gating-Geräusch,
    \item außer Kontrolle geratene Erregung,
    \item unzureichende Refraktärrückgewinnung,
    \item Verlust der Kopplungsstabilität,
    \item unkontrollierte Depolarisation.
\end{itemize}

Fehler, wenn:
\[
T(K_5) > \Theta_{\mathrm{volt-max}}.
\]

% ================================================================
\subsubsection{Energie \texorpdfstring{$E(K_5)$}{E(K_5)}}

Energiefunktional:
\[
E(K_5)
= E_{\mathrm{chem}}
+ E_{\mathrm{ion}}
+ E_{\mathrm{Pumpe}}
+ E_{\mathrm{exc}}
+ E_{\mathrm{Rauschen}}
+ E_{\mathrm{Kopplung}}.
\]

Komponenten:
\begin{itemize}
    \item chemische Energie zur Aufrechterhaltung von Gradienten,
    \item Ionenstromenergie,
    \item pumpengetriebener Energieverbrauch,
    \item Anregungs- und Gating-Energie,
    \item Lärmschutzenergie,
    \item Kopplungsenergie in Netzwerken mit mehreren Einheiten.
\end{itemize}

% ================================================================
\subsubsection{Operatoren auf \texorpdfstring{$K_5$}{K_5} (\texorpdfstring{$\Psi$}{\Psi}, \texorpdfstring{$\Phi$}{\Phi}, \texorpdfstring{$\Lambda$}{\Lambda}, \texorpdfstring{$U$}{U}, \texorpdfstring{$\Chi$}{\Chi})}

\begin{itemize}
    \item $\Psi_{5\to6}$ – Entstehung kognitiver Modellierungsachsen und
          Umwandlung von Anregungszyklen in Attraktoren,
    \item $\Phi$ – Entwicklung der Kanalkinetik und Leitfähigkeiten,
    \item $\Lambda$ – strukturelle Kopplung der Anregung mit der Chemikalie
          Homöostase,
    \item $U$ – Vereinigung ionischer Spezies zu integrierten elektrischen
          Dynamik,
    \item $\Chi$ – Verzweigung in verschiedene Erregbarkeitsregime (Burst,
          tonisch, phasisch).
\end{itemize}

% ================================================================
\subsubsection{Prozesse auf \texorpdfstring{$K_5$}{K_5}}

Typische Prozesse:
\begin{itemize}
    \item Spike-Initiierung und Wiederherstellung,
    \item Gating-Kinetik-Übergänge,
    \item Schwankungen um die Erregungsschwelle,
    \item Ausbreitung von Fronten über $G$,
    \item Feuerfestigkeitsgesteuerte Ordnung der Aktivität,
    \item-Synchronisation in gekoppelten Einheiten,
    \item Entstehung attraktorartiger stabiler Feuermuster.
\end{itemize}

% ================================================================
\subsubsection{Vorhersagen für \texorpdfstring{$K_5$}{K_5}}

\begin{enumerate}
    \item Erregbarkeit erfordert streng kontrollierte Lärmregime.
    \item Spitzenartige Ereignisse treten spontan auf, sobald $\Theta_{\mathrm{exc}}$
          ist gekreuzt.
    \item Die globale Depolarisierung führt zum Strukturkollaps.
    \item Vernetzte $K_5$-Systeme bilden Attraktoren, die $K_6$ vorwegnehmen
          kognitive Modelle.
    \item Kanalvielfalt erweitert lebensfähige Regionen von $\Omega(K_5)$.
\end{enumerate}

% ================================================================
\subsubsection{Experimente für \texorpdfstring{$K_5$}{K_5}}

Empirische Analoga:
\begin{itemize}
    \item künstliche Lipidvesikel mit rekonstituierten Ionenkanälen,
    \item Patch-Clamp-Experimente zur Messung von Gating-Rauschen,
    \item Spannungsklemmencharakterisierung von $\Delta V$-Schwellenwerten,
    \item Multivesikel-Kopplungs- und Synchronisationstests,
    \item Protoneuronenmodelle mit minimaler Erregbarkeit.
\end{itemize}

Diese Tests prüfen direkt $\Theta_{\mathrm{exc}}$, $\Theta_{\mathrm{noise}}$,
$\Theta_{\mathrm{channel-open}}$ usw.

% ================================================================
\subsubsection{Zusammenbruch und Tod von \texorpdfstring{$K_5$}{K_5}}

Tritt auf, wenn:
\[
\Omega(K_5) = \varnothing.
\]

Mechanismen:
\begin{itemize}
    \item unkontrollierte Depolarisation,
    \item-Kanalfehlfunktion oder Gating-Zusammenbruch,
    \item Ionengradientendissipation,
    \item Lärmbedingte Ausfälle durch Durchgehen,
    \item Unfähigkeit, sich von der Erregung zu erholen.
\end{itemize}

Der Kollaps verhindert die Entstehung von $K_6$.

% ================================================================
\subsubsection{Falschbarkeit von \texorpdfstring{$K_5$}{K_5}}

Wichtige überprüfbare Vorhersagen:
\begin{enumerate}
    \item Erregbares Verhalten kann ohne feuerfestes Material nicht aufrechterhalten werden
          Zyklen.
    \item Lärmschutzgrenzwerte sind real und messbar.
    \item Ionenkanalverteilungen bestimmen Lebensfähigkeitsgrenzen.
    \item Stabile Attraktoren in $K_5$ müssen existieren, bevor sich $K_6$ bilden kann.
    \item Spike-Ausbreitung erfordert $\Theta_{\mathrm{front}}$.
\end{enumerate}

Das Scheitern dieser Vorhersagen widerlegt das Modell von $K_5$.

% ================================================================
\subsubsection{Verzweigung / ontologische Position von \texorpdfstring{$K_5$}{K_5}}

Verzweigungstypen:
\begin{itemize}
    \item tonisch-phasische-Burst-Erregbarkeitsregime,
    \item Mehrkanal- vs. Einkanalsysteme,
    \item schwach vs. stark gekoppelte Einheiten,
    \item verteilte vs. lokalisierte Anregungsmuster.
\end{itemize}

Ontologischer Ort:
\[
K_4 \to K_5 \to K_6.
\]

$K_5$ ist die Brücke zwischen biologischen Protozellen und kognitiven Systemen.

% ================================================================
\subsubsection{Beziehung zu M-Räumen}

$M_5$ gibt an:
\begin{itemize}
    \item erlaubte ionische Spezies und Ladungsträger,
    \item zulässige Leitfähigkeits- und Potentialbereiche,
    \item Rauschregime, die mit der Erregbarkeit kompatibel sind,
    \item Umgebungsbedingungen, die die Kanalstabilität ermöglichen,
    \item erlaubte Kopplungsstrukturen.
\end{itemize}

Die Kompatibilität von $K_5$ mit $M_5$ bestimmt die Existenz von
elektrisch erregbare Systeme.
